\section{Excluding small covers at affine pencils}\label{sec:pencils}

The deterministic reduction leaves a specific obstruction at an affine
point $p$.  If $b$ selected split lines pass through $p$, a hypothetical
cover has already spent $q+1-b$ of its labels on the other directions.
Only $b-1$ labels remain to cover the $4b$ designated points, beyond those
covered by partner lines of constituents through $p$.  For $b\ge3$, the
local condition in Section~\ref{sec:deterministic-construction} excludes
precisely this possibility.  We prove that the selection procedure returns
a collection violating this condition with probability $o(1)$.

The input from selection is an upper bound for prescribed choices,
conditional on fixed offsets.  We first apply that bound and then average
geometric events under the original independent offset law.  In
particular, successful return remains part of the event being estimated;
we do not condition the offset law on successful return.

Throughout this section the offsets $M_d^0,M_d^1$ have the independent,
with-replacement law used in the construction.  Let $\mathcal A$ be the
event that the selection procedure returns a compatible collection.  The
conditional estimate \eqref{sel:conditional-offset-cylinder}, obtained
from Proposition~\ref{sel:selection}, says that for any fixed offset
configuration $M$ and any $s$ prescribed tuple choices in distinct slots,
\begin{equation}\label{eq:pencil-cylinder}
 \mathbb P(\mathcal A\text{ and these choices}\mid M)
       \le (K/q)^s,
\end{equation}
where $K$ is the fixed constant in that estimate.  The prescriptions may
depend on $M$, provided they are fixed once $M$ is fixed.  An unavailable
tuple has probability zero, so they may range over all generator and bit
labels.  On $\mathcal A$, $M_d^0\ne M_d^1$ for every direction $d$, no selected
split line is a constituent, and constituent multiplicity at every point
is at most $B=\lceil\log q\rceil$.

\begin{proposition}\label{prop:pencil-local}
In the independent offset experiment with the selection procedure above,
\[
 \mathbb P\bigl(\mathcal A\text{ and the local condition fails at some }
 p\in A\bigr)=o(1).
\]
Here the local condition, required whenever $b(p)\ge3$, is the prohibition
on covers of $F(p)$ by $Z(p)$ and at most $b(p)-1$ additional sets stated in
Section~\ref{sec:deterministic-construction}.
\end{proposition}

The estimate at a fixed point must save three powers of $q$, since we
will sum over $q^2$ affine points.  We obtain those savings from two
sources: repeated generator directions and additional incidences forced
by a cover.  Three generator reuses already cost $O(q^{-3})$.  With zero,
one, or two reuses, we will extract respectively three, two, or one
further incidence costs.  The resulting bound at a fixed point is
$O(B^3/q^3)$, apart from the separately controlled event of a large pencil.

\subsection{Generator reuse in a selected pencil}

Fix $p\in A$.  Write $b$ for the number of selected split lines through
$p$.  Compatibility implies that their target directions are distinct and,
when $b\ge2$, that their $4b$ designated points are distinct and different
from $p$.  Let $\Gamma$ be the union of their generator directions and let
$t_d$ be the number of tuples using $d\in\Gamma$.  Define the reuse count
\[
 \kappa=4b-|\Gamma|=\sum_{d\in\Gamma}(t_d-1).
\]
Thus $\kappa$ counts generator occurrences after their first use.
Compatibility says that two tuples share at most one generator.

We first remove large pencils, so the later enumeration only has to handle
$3\le b\le B$.  The number of selected lines through any point is at most
$B$ except on an event of probability $o(1)$, jointly with $\mathcal A$.
Indeed, for a
specified slot there is at most one retained candidate on the line of its
direction through $p$.  Conditional on $M$, the cylinder estimate and a
union bound give
\[
 \mathbb P(\mathcal A,\ b(p)\ge B\mid M)
 \le \binom{2(q+1)}B(K/q)^B
 \le \left(\frac{C}{B}\right)^B.
\]
The last expression, multiplied by $q^2$, tends to zero.

\begin{lemma}\label{lem:pencil-reuse}
For a fixed affine point $p$, with $\kappa$ the reuse count of the
selected tuples through $p$,
\[
 \mathbb P(\mathcal A,\ b(p)\ge2,\ \kappa\ge3)=O(q^{-3}).
\]
\end{lemma}

\begin{proof}
We extract a bounded collection of tuples witnessing three reuses, then
sum its selection and offset probabilities.  Form a graph whose vertices
are the tuples and join two vertices when their tuples share a generator.
If the number of vertices minus the number of connected components is at
least three, this graph
contains connected pieces with a total of three noninitial vertices and
at most six vertices altogether.  To see this, allocate three units among
the quantities $|C|-1$ for its components, and in each used component grow
a connected set with one more vertex than its allocation.  Order each
piece from a root so that every subsequent vertex has a generator already
used in the piece.

If the number of vertices minus the number of components is at most two,
then $\kappa\ge3$ can occur only in a component of three vertices, with a
different generator on each of the three pairwise overlaps.  Two components
of size two contribute at most two reuses.  In a component of size three,
one generator common to all three contributes exactly two reuses and
precludes any additional overlap.  Otherwise three reuses require the
three distinct pairwise overlaps.  In an ordering of this triangle, the
second tuple has one old generator and the third has two.

In either case the witness has bounded size.  To bound the probability
that one is selected, we bound the expected number of selected witnesses.
A prescription specifies its slots (and hence its target directions),
its generator directions, and its bits.
For a length-$l$ prescription $\pi$, let $T_\pi$ be the offset event that
its tuple lines pass through $p$ and that its positions are distinct and
external to $p$.  By conditioning on $M$ first, its contribution is at most
\begin{equation}\label{eq:pencil-weight}
 (K/q)^l\mathbb P_M(T_\pi).
\end{equation}
Here $\mathbb P_M$ is the unconditioned independent offset law.  Only
after obtaining this bound do we relax restrictions such as distinctness
of slots in the sum of its right side.

We now estimate the sum of \eqref{eq:pencil-weight} by processing the
witness in the order just chosen.  Expose all offsets of a tuple's
generators when that tuple is processed.  Call a generator old if its
direction has already been exposed.  There are a bounded number of
exposed directions in a witness.  If the next tuple has $j$ old generators,
there are $O(q^{4-j})$ choices for its
generator and bit labels.  A position with two old adjacent generators
is fixed; all other positions are independent and uniform either in the
plane or on an exposed constituent.  Their coordinates identifying a
line of the target direction are independent uniform elements of
$\mathbb F_q$.  The maximum number $f$ of fixed positions for
$j=0,1,2,3,4$ is respectively $0,0,1,2,4$.

If no position is fixed, there are $O(q)$ target directions and the
alignment probability with the prescribed line $p+g$ is $q^{-4}$.
If a position is fixed, it must be external to $p$ and therefore determines
$g$; the other $4-f$ positions cost $q^{-(4-f)}$.  Including the two slot
choices and one factor $K/q$ from \eqref{eq:pencil-weight}, the conditional
extension bounds are
\begin{equation}\label{eq:pencil-extension}
\begin{array}{c|ccccc}
j&0&1&2&3&4\\ \hline
\text{extension bound}&O(1)&O(q^{-1})&O(q^{-2})&O(q^{-2})&O(q^{-1}).
\end{array}
\end{equation}
For $j=2$ both possibilities $f=0$ and $f=1$ give $O(q^{-2})$.
The estimates are uniform in previously exposed offsets.  An impossible
fixed-position configuration contributes zero.

Each of the three noninitial steps in the first type of witness saves at
least $q^{-1}$.  The triangle saves $q^{-1}$ and $q^{-2}$ at its second
and third steps.  Summing the finitely many witness-order types and
iterating \eqref{eq:pencil-extension} proves the lemma.  No bound on $b$
was needed.
\end{proof}

\subsection{The aligned law for at most two reuses}

Lemma~\ref{lem:pencil-reuse} handles all pencils with at least three
reuses, regardless of their size.  After also removing the large-pencil
event, it remains to consider $3\le b\le B$ and
$\kappa\in\{0,1,2\}$.  We first describe the offset law for a prescribed
list whose tuple points lie on lines through $p$.  This law will supply
the independent variables and prescribed projections used in the
remaining incidence estimates.

We enumerate ordered lists of $b$ prescribed tuples through $p$.  A pattern
of equality among their $4b$ generator entries with $\kappa$ reuses has
$4b-\kappa$ generator classes.  There are only $b^{O(1)}$ such patterns:
choose the at most two entries that repeat a previous class and their
earlier representatives.  We retain only patterns compatible with at most
one common generator per tuple pair.

For a fixed pattern there are at most
\begin{equation}\label{eq:pencil-directions}
 (q+1)^{5b-\kappa}
\end{equation}
assignments of target and generator directions: there are $b$ target
entries and $4b-\kappa$ generator classes.  Call an assignment
\emph{eligible} if its targets are distinct, its distinct generator
classes have distinct directions, and no tuple has its own target as a
generator.  Bits and slot choices contribute a factor $16^b2^b=32^b$.

For these patterns we may process the tuples as roots and children, each
child sharing exactly one generator with an earlier root and having its
other three generators new.  With one reuse there is one linked pair.
With two reuses there are two disjoint linked pairs, a chain of three
tuples, or one generator direction shared by three tuples.  For a chain
take the middle tuple as root; for a direction shared by three tuples
take any of the three as root.  Other tuples are roots.  This processing
order depends only on the equality
pattern, not on the direction assignment or offsets.

Write $g_i$ for the target direction of tuple $i$ and put $L_i=p+g_i$.
We call the list \emph{aligned} if all four positions of each tuple $i$
lie on $L_i$.  This event imposes neither distinctness of the positions
nor any other validity condition.  For every eligible direction
assignment,
\begin{equation}\label{eq:pencil-alignment}
 \mathbb P_M(\text{alignment})=q^{-4b}.
\end{equation}
Conditional only on alignment, the offsets and positions can be generated
as follows; we refer to this distribution as the \emph{aligned law}.
\begin{enumerate}[label=(\roman*)]
\item At a root, choose four independent uniform points of $L_i$ and set
the generator offsets to the prescribed lines through their endpoints.
\item At a child, the two positions involving its shared generator are
the intersections of that generator's already fixed constituents with
$L_i$.  Choose the other two positions independently and uniformly on
$L_i$, and set the new generator offsets accordingly.
\item All offsets in directions outside $\Gamma$ retain their original
independent law.
\end{enumerate}
For a root, before alignment its four positions are independent uniform
plane points, giving the factor $q^{-4}$.  At a child, the two positions
on the shared generator's constituents each involve one fresh offset.
Each lands at its required intersection with probability $1/q$.  The
other two positions are independent uniform plane points, each costing
$1/q$ to lie on $L_i$.  These use disjoint fresh offset variables.  The
factor $q^{-4}$ is independent of the earlier values, so later alignment
does not bias an earlier generation.  This proves both
\eqref{eq:pencil-alignment} and the stated law.  A target of one tuple may
equal a generator of another; that does not affect this argument.

The aligned law deliberately leaves the remaining geometric restrictions
as events to be tested.  Call a list \emph{admissible} if it is aligned,
all its positions are distinct and external to $p$, none of its $L_i$ is
a constituent, and the constituent multiplicity is at most $B$ at every
point.  Every list selected on $\mathcal A$ through $p$, with $b\ge3$,
satisfies these necessary conditions.

To express a violation of the local condition for a prescribed list,
let $F$ be the union of its designated points and define $Z(p)$ by
\eqref{eq:det-partner-set}, with $F$ in place of $F(p)$; throughout
these counts $Z(p)$ refers to this prescribed list.
A \emph{bad cover} of this list covers $F$ by $Z(p)$ and at most $b-1$
additional sets: sets of at most two points, affine lines outside
$\{L_1,\ldots,L_b\}\cup\mathcal M$, or merged pairs, with no merged
pair used more than once. An admissible bad cover means that both
this cover and the preceding admissibility conditions occur.
In the following estimates we bound a bad cover jointly with
admissibility, rather than conditioning the point variables on it.

For a used direction $d\in\Gamma$, its $2t_d$ prescribed endpoint
indices will be called its \emph{purposeful endpoints}.  Neither
constituent of $d$ passes through $p$ on an admissible list.  Indeed, it
contains a purposeful point external to $p$ on a target line through $p$,
and its direction differs from that target. Consequently the set $Z(p)$
of points already covered by partner lines uses only directions outside
$\Gamma$; we call this precoverage.

\subsection{Additional incidence costs for zero or one reuse}

For zero or one reuse, we can estimate each eligible direction assignment
separately.  Even after allowing the purposeful endpoints covered by used
merged pairs, a cover by $b-1$ additional sets must create further
incidences.  We will extract $3-\kappa$ of these incidence conditions
whose variable dependencies permit successive integration.

\begin{lemma}\label{lem:pencil-small-reuse}
Fix $p\in A$, $3\le b\le B$, an equality pattern with $\kappa=0$ or $1$,
an eligible direction assignment, and the bits.  Under the aligned law,
the probability of an admissible bad cover is at most
\[
 b^{O(1)}(2B/q)^{3-\kappa}.
\]
The polynomial exponent and implicit constant are absolute.
\end{lemma}

\begin{proof}
When $\kappa=0$, all $4b$ positions are independent uniform points of
their target lines.  When $\kappa=1$, two positions of the child are
projections of two parent positions along the shared generator direction.
Call these two child indices dependent.  All nondependent positions are
independent uniform points of their target lines, and each has at most
one dependent image.  We fix the offsets in unused directions, restricting
to configurations whose unused constituent multiplicity is at most $B$
everywhere.  These offsets are independent of the point variables under
the aligned law.  Configurations violating this multiplicity bound
contribute no admissible occurrence.

We first turn a cover into incidence conditions.  Given an admissible
bad cover, assign each of the $4b$ positions to precoverage or to one of
at most $b-1$ additional covering sets.  Mark certain assigned indices as
\emph{surplus}, according to the following rule:
\begin{itemize}
\item Mark every position assigned to precoverage.
\item For an ordinary line or an unused merged pair with at least three
assigned positions, retain two assigned indices as anchors and mark all
others.  A set of size at most two needs no marks.
\item For a used merged pair of direction $d$ with $n'$ assigned
positions, mark $\max\{0,n'-2t_d\}$ nonpurposeful indices.  There are
enough such indices because there are only $2t_d$ purposeful endpoints.
\end{itemize}
Every additional set retains at most two unmarked indices, except that a
used merged pair may retain an additional $2(t_d-1)$.  Each used pair is
allowed at most once.  Thus the number of surplus indices is at least
\[
 4b-2(b-1)-2\sum_{d\in\Gamma}(t_d-1)
       =2b+2-2\kappa.
\]
Discard any dependent indices from the surplus and denote the remainder
by $I$.  Then
\begin{equation}\label{eq:pencil-surplus}
 |I|\ge 2b+2\quad(\kappa=0),\qquad
 |I|\ge 2b-2\quad(\kappa=1).
\end{equation}

For each $v\in I$, its covering assignment gives an incidence condition
on its independent point variable $x_v$, using at most two anchor
variables:
\begin{enumerate}[label=(\alph*)]
\item In precoverage, $x_v$ lies on a partner of an unused constituent
through $p$.
\item On an ordinary line, $x_v$ lies on the line through its two stored
anchors, which are distinct points.
\item On an unused merged pair, choose one stored anchor whose underlying
variable is not $x_v$.  Then $x_v$ lies on either constituent of some
unused merged pair containing that anchor.
\item On a used merged pair, $x_v$ lies on the line of that generator
direction through one of its purposeful endpoints.  Choose an endpoint
on the constituent actually covering $x_v$.
\end{enumerate}
In every case the candidate line containing $x_v$ differs from its own
target line.  Replace a dependent anchor by its underlying parent
variable when recording these dependencies.

No condition needs its own variable as an anchor.  Only the case
$\kappa=1$ requires explanation.  For an ordinary line, $x_v$ and its
dependent image lie on a shared-generator constituent, which is forbidden
as an ordinary covering line.  For an unused pair, the two stored anchors
are different from the surplus index, and at most one can be its
dependent image; choose the other.  For a used pair, the surplus index is
nonpurposeful for that generator.  If its chosen endpoint were its
dependent image, the line joining the two distinct points would have the
shared generator direction.  The surplus index would then itself be
purposeful for that direction, a contradiction.

The surplus count alone does not allow us to multiply incidence
probabilities: one condition may depend on another condition's variable.
We therefore extract a collection with acyclic dependencies.  Draw a
directed edge from $v\in I$ to each underlying anchor variable that also
belongs to $I$.  This directed graph has no loops and outdegree
at most two.  It has an acyclic induced set of size at least
$\lceil |I|/3\rceil$: process the vertices in any order, assigning each
one of three colours not used by its already processed outneighbors.  A
monochromatic directed cycle is impossible, since its last processed
vertex has an earlier outneighbor on the cycle.  Each colour class is
therefore acyclic.  By \eqref{eq:pencil-surplus} and $b\ge3$, we can retain
exactly $3-\kappa$ incidence conditions inducing an acyclic graph.

Only the retained conditions now need to be counted.  There are
$b^{O(1)}$ possible specifications of their indices, condition types,
anchor indices, and, when needed, used-generator direction, because at
most three conditions remain.  In case (c) the unused direction is left
unspecified.  Thus we take a union over these specifications without
enumerating the full cover.  By the extraction, it suffices to include
specifications with no self-dependency and with an acyclic dependency
graph.

Fix such a specification and all independent variables outside its
retained indices.  The acyclic graph gives an integration order: integrate
the remaining independent uniform variables by
successively removing a vertex of indegree zero.  No other remaining
condition depends on the variable removed.  With the other variables
fixed, its condition is either impossible or asks for membership in at
most $2B$ lines different from its target line.  In case (a) there are at
most $B$ relevant constituents through $p$; in case (c) the fixed anchor
belongs to at most $B$ unused constituents, giving at most $2B$ candidate
constituents of the corresponding pairs.  Cases (b) and (d) give one
line.  Each candidate line meets the target line in at most one point.
The conditional probability is therefore at most $2B/q$ at every removal.

For this integration we retain only the indicated incidence conditions,
their line exclusions, and the distinctness of any two ordinary-line
anchors needed to define their line.  We discard other cover, distinctness, and
admissibility conditions.  Their role was to extract a specification with
no self-dependency; retaining them as conditioning would change the law
of the variables being integrated.
The specified event costs at most $(2B/q)^{3-\kappa}$.  Sum over the
polynomially many specifications and average over the fixed unused offsets.
\end{proof}

\subsection{One additional incidence when there are two reuses}

Two reuses already contribute two powers of $q^{-1}$ in the final
prescription count.  It remains to gain one further power from the bad
cover.  Here we average over direction assignments as well as aligned
positions.  Each incidence estimate will reveal a direction after the
geometric data prescribing its required value are known.

\begin{lemma}\label{lem:pencil-two-reuses}
Fix $p\in A$, $3\le b\le B$, an equality pattern with $\kappa=2$, and
the bits.  Summed over eligible direction assignments, the probabilities
of alignment and an admissible bad cover are at most
\[
 (q+1)^{5b-2}q^{-4b}\,O(b^3/q),\qquad 3\le b\le B.
\]
\end{lemma}

\begin{proof}
For the direction count, regard each target and each generator equality
class as one \emph{token}; there are $5b-2$ tokens.  We first estimate the
mean probability of bad geometry over assignments in which all tokens
receive distinct directions.  Equivalently, assign them by a uniform
injection into $\mathcal G$.  At the end we bound the remaining eligible
assignments, in which a target direction equals a generator direction of
another tuple.  Under the aligned law an admissible bad cover requires at
least one of the following four events:
\begin{enumerate}[label=(\roman*)]
\item An unused direction has one constituent through $p$ and the other
through a designated point.
\item A used merged pair meets a point of a tuple not using that generator.
\item Three designated points lie on a line different from every target
line and from every constituent of a used direction.
\item Three noncollinear designated points lie in one unused merged pair.
\end{enumerate}
Indeed, if (i) fails there is no precoverage.  If (ii) fails, a used pair
covers only its $2t_d$ purposeful endpoints.  Ordinary allowed lines cover
at most two points when (iii) fails.  If an unused pair covers three points,
they either witness (iv) or are collinear.  In the latter case two lie on
one constituent, so their common line is that unused constituent.  By
admissibility it is not a target line, and it cannot be a constituent of
a used direction; thus it witnesses (iii).  Without these events the
additional $b-1$ sets cover at most
\[
 2(b-1)+2\kappa=2b+2<4b
\]
points, a contradiction.

Conditionally on directions and designated points, unused offsets are
independent.  Event (i) has probability $O(b/q)$: for a specified point,
direction, and orientation, the two offsets have prescribed values and
cost $q^{-2}$, and there are $O(bq)$ choices.  Event (iv) costs
$O(b^3/q^2)$: choose the three noncollinear points and two that lie on one
constituent.  Their joining line determines the direction and one offset;
the parallel line through the third point determines the other offset.
If that direction is unused this costs $q^{-2}$, with two possible offset
orientations.

The remaining events (ii) and (iii) involve used constituents and may
involve projected positions.  Their estimates use the randomness of the
direction assignment.  To keep the relevant conditional laws explicit,
use the following sequential version of the injection average.  Process
roots and children in the order already chosen.  First
reveal a tuple's target token, then generate its fresh positions and any
prescribed projections, and finally reveal its new generator tokens and
set their offsets.  This has the aligned law: positions use the present
target and any previously fixed shared constituents, but not the values
of the new generator directions.  In particular, a projected position
need not be an independent uniform point.  When the tested line is fixed
before that position's tuple is processed, its incidence is tested at the
reveal of the tuple's target token; the opposite timing is handled below
by revealing the later generator token.  The alignment factor $q^{-4b}$
is the same for every injection.  At every stage
unrevealed tokens are uniformly
assigned without replacement among unused directions.  In particular a
specific required value determined before a token is revealed costs at
most
\begin{equation}\label{eq:pencil-token}
 (q+1-5b)^{-1}=O(q^{-1}).
\end{equation}

Consider (ii) for a specified used constituent and a specified point of
a tuple not using that generator.  If the constituent was fixed before
that tuple is processed, a fresh position hits it with probability $1/q$.
A dependent position is the projection of an already fixed shared
constituent onto the new target line.  To hit the tested constituent,
the target line must pass through the intersection of the two fixed
constituents.  An intersection at $p$ cannot produce an admissible external
point.  Any other intersection prescribes a target direction and costs
$O(q^{-1})$ by \eqref{eq:pencil-token}.

In the opposite timing, the tested point is known before the first tuple
using the tested direction is processed.  The purposeful endpoint defining
the tested constituent is generated before that new generator token is
revealed.  Admissibility makes the endpoint different from the earlier
point.  Their joining direction is therefore prescribed at the reveal of
the generator token, again costing $O(q^{-1})$.  Summing over $O(b^2)$
choices bounds (ii), jointly with the indicated admissibility conditions,
by $O(b^2/q)$.

For (iii), the three distinct points come from three distinct tuples:
two on one target line would force the triple line to be that target.
The points in the first two processed tuples determine a line $h$.
At the last tuple, a fresh position hits $h$ with probability at most
$1/q$, keeping the necessary exclusion $h\ne L_i$.  If the last position
is a projection along a shared constituent $h_s$, the event requires
$h\ne h_s$.  Parallel lines give no incidence.  Otherwise their
intersection must differ from $p$ and prescribes the new target direction,
giving $O(q^{-1})$ by \eqref{eq:pencil-token}.  Summing over triples gives
$O(b^3/q)$.  Every constraint in this argument is determined before the
random value used to enforce it is revealed.  Other conditions are
discarded in taking these upper bounds.

Combining the four events, the mean bad-geometry probability over
injective token assignments is $O(b^3/q)$.  Multiplying by the common
alignment probability $q^{-4b}$ and by the number of injections gives
the required contribution from these assignments.  The only eligible
noninjective assignments have an equality between a target token and a
generator token.  There are at most
\[
 b(4b-2)(q+1)^{5b-3}
\]
such assignments, by selecting one cross-equality and ignoring remaining
restrictions.  For each eligible assignment the alignment probability is
still exactly $q^{-4b}$; use the trivial upper bound one for subsequent
bad geometry.  Their total is absorbed by the claimed bound.
\end{proof}

\subsection{Summing over all selected pencils}

We now combine the selection estimate with the geometric estimates to
prove Proposition~\ref{prop:pencil-local}.  The former costs one factor
$K/q$ per selected tuple; the latter supplies the alignment probability
and the additional incidence costs.  For fixed $p,b,\kappa$ with
$3\le b\le B$ and $\kappa\le2$, count ordered prescriptions that occur
as selections on $\mathcal A$ and have an admissible bad cover using the
precoverage defined by that list.  An actual violation with exactly these
parameters yields all $b!$ orders of
its selected tuples.  The count need not require that there are no further
selected lines through $p$: dropping that restriction only increases it.

For each prescription, condition on $M$ first to obtain the factor
$(K/q)^b$ from \eqref{eq:pencil-cylinder}.  All necessary geometric tests
then depend only on the independent offset model and that prescription.
The patterns, slots, bits, direction assignments, and alignment contribute,
after division by $b!$, at most
\begin{equation}\label{eq:pencil-common-factor}
 \frac{b^{O(1)}}{b!}\,32^b(q+1)^{5b-\kappa}
       (K/q)^b q^{-4b}
 \le \frac{b^{O(1)}C^b}{b!}q^{-\kappa}.
\end{equation}
Here $C$ and the polynomial exponent are independent of $b$ and $q$;
we use $b\le\lceil\log q\rceil$.  For $\kappa=0,1$, multiply by the
additional factor from Lemma~\ref{lem:pencil-small-reuse}.  For $\kappa=2$,
Lemma~\ref{lem:pencil-two-reuses} provides the factor $O(b^3/q)$ after
summing direction assignments.  Hence the total probability at $p$ for
these values of $b,\kappa$ is
\[
 O(B^3/q^3),
\]
because $\sum_{b\ge3} b^a C^b/b!<\infty$ for every fixed $a,C$.
Lemma~\ref{lem:pencil-reuse} adds $O(q^{-3})$ for all cases with
$\kappa\ge3$.  Summing over the $q^2$ affine points gives
$O(B^3/q)=o(1)$.  The earlier large-$b$ estimate is also $o(1)$.
This proves Proposition~\ref{prop:pencil-local}: a violation of the local
condition occurs with probability $o(1)$, jointly with successful return
of the selection procedure.
